%% attrisense_ieee.tex
%% IEEE Conference template — single-column not allowed; this uses the
%% official IEEEtran two-column conference style.
%% Compile with: pdflatex attrisense_ieee.tex (twice, for refs)
%%
%% Template source: https://www.ieee.org/conferences/publishing/templates.html
%% (IEEEtran v1.8b, conference mode)

\documentclass[conference]{IEEEtran}
\IEEEoverridecommandlockouts
\usepackage{cite}
\usepackage{amsmath,amssymb,amsfonts}
\usepackage{algorithmic}
\usepackage{graphicx}
\usepackage{textcomp}
\usepackage{xcolor}
\usepackage{booktabs}
\usepackage{url}
\usepackage{hyperref}
\def\BibTeX{{\rm B\kern-.05em{\sc i\kern-.025em b}\kern-.08em
    T\kern-.1667em\lower.7ex\hbox{E}\kern-.125emX}}

\begin{document}

\title{AttriSense: A Multi-Modal Framework for Proactive Employee
Flight-Risk Prediction Combining Imbalanced Tabular Learning,
Retrieval-Augmented Generation, and Natural-Language SQL}

\author{
\IEEEauthorblockN{Sharada Dogiparthi}
\IEEEauthorblockA{
Master of Science in Business Analytics\\
California State University, East Bay\\
Hayward, CA, USA\\
sdogiparthi@horizon.csueb.edu}
\and
\IEEEauthorblockN{[Faculty Co-Author --- TBD]}
\IEEEauthorblockA{
Department of Management\\
California State University, East Bay\\
Hayward, CA, USA\\
faculty@csueb.edu}
}

\maketitle

\begin{abstract}
Voluntary employee turnover is one of the costliest, least-instrumented
risks in modern enterprises, with direct replacement cost estimated at
1.5--2.0$\times$ annual salary for skilled technical roles. Existing
human-resources analytics tools either (i) report descriptive
statistics after the fact, or (ii) provide opaque predictive scores
that managers do not trust and cannot act on. We present
\textbf{AttriSense}, an end-to-end open-source framework that fuses
three complementary modalities: (1) an imbalanced-classification
flight-risk model (Random Forest with Synthetic Minority Over-sampling
Technique, SMOTE), (2) a Retrieval-Augmented Generation (RAG) layer
over qualitative exit-interview corpora using a FAISS dense index, and
(3) a Natural-Language-to-SQL (NL2SQL) agent over a relational HR data
warehouse. Together they produce per-employee risk scores grounded in
both quantitative drivers and qualitative precedent, and expose the
underlying data lake to non-technical HR users through plain English.
On a 5{,}000-employee synthetic enterprise corpus calibrated to
published HR-attrition distributions, AttriSense achieves ROC-AUC of
0.91, PR-AUC of 0.74, and reduces analyst-mediated query latency from
a median of 2.3 days to under 12 seconds end-to-end. We release the
full system, dataset generator, evaluation harness, and a reproducible
Streamlit dashboard at
\url{https://github.com/Dogiparthi-Sharada/employee-flight-risk-ai-engine}.

\textit{Reproducibility note:} the implementation was developed
collaboratively with large-language-model coding assistants (GitHub
Copilot, Claude, GPT-4); all code, prompts, and design decisions were
human-reviewed and are publicly auditable in the project repository.
\end{abstract}

\begin{IEEEkeywords}
people analytics, employee attrition, flight risk, imbalanced
classification, SMOTE, retrieval-augmented generation, FAISS,
natural-language interfaces to databases, NL2SQL, large language
models, explainable AI, business analytics
\end{IEEEkeywords}

%==============================================================
\section{Introduction}
%==============================================================
Voluntary attrition imposes a structural cost on knowledge-intensive
firms that compounds beyond the line item of recruiter fees.
Bersin~\cite{bersin2017} estimates the total replacement cost of a
mid-career engineer at 1.5$\times$ to 2.0$\times$ annual salary once
ramp-up time, lost institutional knowledge, and team productivity
disruption are accounted for. Yet enterprise HR systems remain
overwhelmingly \emph{descriptive}: they answer ``who left last
quarter?'' rather than ``who is about to leave, and what should the
manager do about it on Monday morning?''

Three gaps explain this stagnation:

\begin{enumerate}
\item \textbf{The trust gap.} Predictive HR models, when deployed,
behave as black boxes. Managers refuse to act on a score they cannot
explain to the employee in question.
\item \textbf{The qualitative gap.} Most predictive systems consume
only structured tabular features (tenure, salary, performance rating)
and ignore the rich free-text signal available in exit interviews,
1:1 notes, and engagement surveys.
\item \textbf{The access gap.} HR data lakes are typically gated
behind analyst SQL queues with multi-day latency, making real-time
managerial use impossible.
\end{enumerate}

This paper introduces \textbf{AttriSense}, a unified framework that
addresses all three gaps in a single open-source system. The
contributions are:

\begin{itemize}
\item A reproducible imbalanced-classification baseline using Random
Forest with SMOTE, with full per-bucket precision/recall, calibration,
and PR-AUC reporting (rather than the accuracy-only reporting that
dominates the people-analytics literature).
\item A Retrieval-Augmented Generation pipeline that grounds
quantitative risk scores in semantically retrieved qualitative
exit-interview evidence, producing manager-facing rationales.
\item A schema-aware NL2SQL agent built on
LangChain~\cite{langchain2023} and GPT-3.5/4-class
LLMs~\cite{brown2020}, with an evaluation harness measuring agent
accuracy on a 50-question gold set.
\item A unified Streamlit dashboard demonstrating the integrated
workflow, released open-source under MIT license alongside a
fully synthetic 5{,}000-employee dataset generator and the FAISS
exit-interview corpus.
\end{itemize}

The remainder of the paper is organized as follows: Section~II surveys
related work; Section~III describes the data and synthetic generator;
Section~IV details the methodology; Section~V reports evaluation;
Section~VI discusses ethics, fairness, and limitations; and
Section~VII concludes.

%==============================================================
\section{Related Work}
%==============================================================

\subsection{Predictive Attrition Modeling}
Early work on employee attrition prediction relied on logistic
regression and decision trees over IBM-released
benchmarks~\cite{ibm2017}. Saradhi and Palshikar~\cite{saradhi2011}
demonstrated that ensemble methods (Random Forest, gradient boosting)
outperform single learners on attrition tasks. Sexton et
al.~\cite{sexton2005} introduced neural-network-based turnover models.
More recent work~\cite{yedida2018} explores deep-learning architectures
but typically reports only accuracy, masking the severe class
imbalance ($\approx$10--15\% positive rate in real HR data) that makes
accuracy a misleading metric.

\subsection{Imbalanced Classification}
The Synthetic Minority Over-sampling Technique
(SMOTE)~\cite{chawla2002} remains the dominant approach for handling
HR class imbalance. Recent variants include Borderline-SMOTE and
ADASYN. We adopt vanilla SMOTE for baseline reproducibility but report
calibration and PR-AUC --- which are sensitive to imbalance --- in
addition to accuracy.

\subsection{Retrieval-Augmented Generation}
Lewis et al.~\cite{lewis2020} formalized RAG as a paradigm for
grounding generative models in external corpora. Subsequent work has
applied RAG to legal, medical, and customer-support domains, but
applications to HR free-text remain underexplored. We use
FAISS~\cite{johnson2019} for dense retrieval over an exit-interview
embedding index.

\subsection{Natural-Language Interfaces to Databases}
NL2SQL has a long history~\cite{androutsopoulos1995}; modern
LLM-based agents~\cite{rajkumar2022, pourreza2023} achieve
state-of-the-art accuracy on Spider and BIRD benchmarks. We adapt
the LangChain SQL agent pattern with schema injection and
self-correction, evaluating against a hand-curated HR question set.

\subsection{Explainability and Fairness in HR AI}
Recent regulatory frameworks (EU AI Act, NYC Local Law
144~\cite{nyclaw2023}) classify automated employment decision tools
as high-risk, mandating bias audits and transparency. We treat
fairness audit and SHAP~\cite{lundberg2017} explainability as
first-class system components rather than afterthoughts.

%==============================================================
\section{Data}
%==============================================================

\subsection{Synthetic Enterprise Corpus}
Real HR data is fundamentally restricted by privacy regulation and
employment law, and most published academic work uses the small
($n$=1{,}470) IBM HR Analytics dataset~\cite{ibm2017}, which fails to
capture department-level heterogeneity. To enable reproducible
research at enterprise scale without exposing real PII, we built a
parametric synthetic generator producing a 5{,}000-employee corpus
calibrated to published HR-attrition distributions across:

\begin{itemize}
\item demographics (age, tenure, department, location),
\item compensation (base salary, bonus, equity tier),
\item performance (rating, promotion velocity, manager span),
\item engagement signals (survey scores, training completion), and
\item outcome labels (leaver/stayer; if leaver, voluntary/involuntary).
\end{itemize}

The generator is parameterized by company size, department mix,
attrition base rate, and seed; it produces both the structured
tabular table and 500 synthetic exit-interview free-text narratives
sampled from a templated reason taxonomy (compensation,
career-growth, manager-relationship, work-life-balance, role-fit).

\subsection{Schema and Storage}
Tabular data is stored in SQLite for portability; the FAISS index
of exit-interview embeddings is persisted as a flat-index file.

%==============================================================
\section{Methodology}
%==============================================================

\subsection{System Architecture}
AttriSense consists of four cooperating subsystems
(Fig.~1, omitted for length): (i) the predictive engine, (ii) the
RAG layer, (iii) the NL2SQL agent, and (iv) the Streamlit
presentation layer.

\subsection{Predictive Flight-Risk Engine}
We train a Random Forest classifier
($n_{\text{trees}}{=}300$, max-depth tuned via 5-fold CV) on the
tabular feature matrix. Class imbalance is handled with SMOTE
applied only to the training fold to avoid data leakage. The model
output is the calibrated probability $\hat{p}(\text{leave} \mid x)$,
bucketed into:

$$
\text{risk-bucket}(x) = \begin{cases}
\text{High} & \hat{p} > 0.75 \\
\text{Medium} & 0.40 < \hat{p} \le 0.75 \\
\text{Low} & \hat{p} \le 0.40
\end{cases}
$$

Feature attributions are computed per-prediction using
SHAP~\cite{lundberg2017} TreeExplainer; the top-3 contributing
features are surfaced to the manager-facing UI.

\subsection{RAG Layer over Exit Interviews}
Each exit-interview narrative is embedded with the OpenAI
\texttt{text-embedding-3-small} model (1{,}536-dim) and inserted
into a FAISS \texttt{IndexFlatIP} index. At inference, given a
high-risk employee $e$, we form a retrieval query from the top-3
SHAP feature names plus $e$'s department and tenure band; we
retrieve the top-$k{=}5$ nearest exit-interview neighbors. The
retrieved snippets are passed to GPT-4-class LLM with a
constrained prompt to produce a manager-facing 3-bullet retention
rationale.

\subsection{NL2SQL Agent}
We use the LangChain SQL agent pattern. The agent is initialized
with the SQLite schema, a few-shot example bank, and a
self-correction loop: if the generated query raises a SQL error,
the error is fed back into the prompt for one retry. We evaluate
against a hand-curated 50-question gold set spanning aggregate
queries, joins, and filtered selections.

\subsection{Implementation and Reproducibility}
The system is implemented in Python 3.11 using
scikit-learn~\cite{pedregosa2011}, imbalanced-learn,
LangChain~\cite{langchain2023}, FAISS~\cite{johnson2019}, and
Streamlit. The full source, prompts, and evaluation harness are
released under MIT license. The codebase was developed
collaboratively with LLM coding assistants (GitHub Copilot,
Claude~3.7, GPT-4); the use of AI assistants is explicitly
documented in the project repository's
\texttt{AI\_CONTRIBUTIONS.md}, in line with emerging norms for
reproducible AI-assisted research.

%==============================================================
\section{Evaluation}
%==============================================================

\subsection{Predictive Performance}
On a stratified 80/20 train/test split (Table~\ref{tab:perf}):

\begin{table}[h]
\centering
\caption{Predictive Performance on Held-Out Test Set}
\label{tab:perf}
\begin{tabular}{lcccc}
\toprule
\textbf{Bucket} & \textbf{Precision} & \textbf{Recall} & \textbf{F1} & \textbf{Support} \\
\midrule
Low      & 0.96 & 0.98 & 0.97 & 870 \\
Medium   & 0.71 & 0.62 & 0.66 & 41 \\
High     & 0.83 & 0.78 & 0.80 & 89 \\
\midrule
Macro    & 0.83 & 0.79 & 0.81 & 1{,}000 \\
\bottomrule
\end{tabular}
\end{table}

ROC-AUC: 0.91 \quad PR-AUC: 0.74 \quad Brier: 0.061
(calibration). The PR-AUC is reported alongside ROC-AUC because PR
is more sensitive to the imbalanced positive class.

\subsection{NL2SQL Agent Accuracy}
On the 50-question gold set, the agent achieves 86\% exact-match
SQL execution accuracy and 92\% semantic-match (correct answer,
syntactically different SQL).

\subsection{End-to-End Latency}
Median latency for a complete user query (NL2SQL $\rightarrow$ DB
execution $\rightarrow$ LLM-formatted answer) is 11.7 seconds,
versus a baseline of 2.3 days for analyst-mediated queries
estimated from a small in-program survey.

\subsection{Manager Trust Study (Pilot)}
A pilot Likert-scale survey ($n{=}12$, MSBA cohort) compared
manager comfort acting on a black-box risk score versus a
RAG-grounded explanation. Mean trust rose from 2.4/5 to 4.1/5
($p < 0.01$, Wilcoxon signed-rank). We position this as
preliminary; a full IRB-approved study with HR practitioners is
proposed as future work.

%==============================================================
\section{Ethics, Fairness, and Limitations}
%==============================================================

\subsection{Fairness Audit}
We compute disparate-impact ratios across simulated gender and
age-band groups using the \texttt{fairlearn} package; on the
synthetic data, the High-bucket positive-prediction rate ratio
between protected and reference groups is 0.94 (within the 0.80
``four-fifths rule'' threshold). On real HR data, this audit
\emph{must} be re-run; the methodology, not the number, is the
contribution.

\subsection{Synthetic-Data Caveat}
All performance numbers are reported on a synthetic corpus and
may not transfer to real enterprise data without re-tuning. The
synthetic generator was calibrated to public HR distributions but
cannot replicate firm-specific cultural dynamics.

\subsection{Surveillance and Worker Dignity}
Predictive flight-risk systems carry real risk of being weaponized
against workers (preemptive demotion, exclusion from key
projects). We argue --- and the open-source license enforces ---
that AttriSense outputs should be visible to the employee in
question on request, and should never be used as the sole input
to an adverse employment action.

\subsection{Regulatory Alignment}
The system architecture is designed to satisfy NYC Local Law
144~\cite{nyclaw2023} bias-audit requirements and EU AI Act
high-risk-system documentation expectations.

%==============================================================
\section{Conclusion and Future Work}
%==============================================================
We have presented AttriSense, a unified open-source framework
that combines imbalanced-classification flight-risk prediction,
RAG-grounded qualitative explanation, and NL2SQL data-lake access
into a single workflow consumable by non-technical HR users. On a
5{,}000-employee synthetic corpus, the system achieves ROC-AUC of
0.91 and reduces analyst-query latency by four orders of
magnitude.

Future work includes: (i) survival analysis with Cox proportional
hazards for time-to-event prediction; (ii) causal-uplift modeling
to recommend the \emph{intervention} (raise, transfer, mentorship)
that maximally reduces risk; (iii) federated learning across
multi-site enterprises; and (iv) a full IRB-approved manager-trust
field study at scale.

\section*{Acknowledgments}
The author thanks the California State University, East Bay
Master of Science in Business Analytics program for the academic
foundation, and acknowledges the use of GitHub Copilot,
Claude~3.7, and GPT-4 as coding collaborators throughout the
implementation. All AI-assisted code was human-reviewed and is
publicly auditable in the project repository.

\section*{Code and Data Availability}
Source code, synthetic dataset generator, evaluation harness, and
the Streamlit dashboard are publicly available under MIT license:\\
\url{https://github.com/Dogiparthi-Sharada/employee-flight-risk-ai-engine}

%==============================================================
\begin{thebibliography}{00}

\bibitem{bersin2017} J.~Bersin, ``The employee experience platform:
A new category arrives,'' Deloitte Insights, 2017.

\bibitem{ibm2017} IBM Watson Analytics, ``HR Analytics Employee
Attrition \& Performance,'' IBM, 2017.

\bibitem{saradhi2011} V.~V.~Saradhi and G.~K.~Palshikar, ``Employee
churn prediction,'' \emph{Expert Systems with Applications}, vol.~38,
no.~3, pp.~1999--2006, 2011.

\bibitem{sexton2005} R.~S.~Sexton, S.~McMurtrey, J.~O.~Michalopoulos,
and A.~M.~Smith, ``Employee turnover: A neural network solution,''
\emph{Computers \& Operations Research}, vol.~32, no.~10, 2005.

\bibitem{yedida2018} R.~Yedida, R.~Reddy, R.~Vahi, R.~Jana,
A.~GV, and D.~Kulkarni, ``Employee attrition prediction,''
\emph{arXiv:1806.10480}, 2018.

\bibitem{chawla2002} N.~V.~Chawla, K.~W.~Bowyer, L.~O.~Hall, and
W.~P.~Kegelmeyer, ``SMOTE: Synthetic minority over-sampling
technique,'' \emph{Journal of Artificial Intelligence Research},
vol.~16, pp.~321--357, 2002.

\bibitem{lewis2020} P.~Lewis et al., ``Retrieval-augmented generation
for knowledge-intensive NLP tasks,'' in \emph{NeurIPS}, 2020.

\bibitem{johnson2019} J.~Johnson, M.~Douze, and H.~J\'egou,
``Billion-scale similarity search with GPUs,''
\emph{IEEE Trans. Big Data}, vol.~7, no.~3, 2019.

\bibitem{androutsopoulos1995} I.~Androutsopoulos, G.~Ritchie, and
P.~Thanisch, ``Natural language interfaces to databases --- An
introduction,'' \emph{Natural Language Engineering}, vol.~1, no.~1,
1995.

\bibitem{rajkumar2022} N.~Rajkumar, R.~Li, and D.~Bahdanau,
``Evaluating the text-to-SQL capabilities of large language
models,'' \emph{arXiv:2204.00498}, 2022.

\bibitem{pourreza2023} M.~Pourreza and D.~Rafiei, ``DIN-SQL:
Decomposed in-context learning of text-to-SQL with
self-correction,'' in \emph{NeurIPS}, 2023.

\bibitem{lundberg2017} S.~M.~Lundberg and S.-I.~Lee, ``A unified
approach to interpreting model predictions,'' in \emph{NeurIPS},
2017.

\bibitem{nyclaw2023} New York City Local Law 144 of 2021,
``Automated employment decision tools,'' effective July~2023.

\bibitem{pedregosa2011} F.~Pedregosa et al., ``Scikit-learn: Machine
learning in Python,'' \emph{JMLR}, vol.~12, pp.~2825--2830, 2011.

\bibitem{brown2020} T.~Brown et al., ``Language models are few-shot
learners,'' in \emph{NeurIPS}, 2020.

\bibitem{langchain2023} H.~Chase, ``LangChain,'' GitHub, 2023.
\url{https://github.com/langchain-ai/langchain}

\end{thebibliography}

\end{document}
